% -------- Subsection 11.3.10 NS projected transport and closure packet --------
% (File: 11.3.10-NSProjectedTransportAndClosure.tex | Volume 11)

\subsection{NS projected transport and closure packet}\label{sub:NSProjectedTransportAndClosurePacket}

This subsection isolates the Navier--Stokes lane-owned theorem packet.
The shared physical-sector selector, the persistent invariant, and the exact/persistent/tail split are treated as imported objects.
The present packet records the continuum coherent-sector route, the projected transport law, the shadow closure theorem, and the exact remaining analytic burdens.

\begin{definition}[Continuum coherent sector and projected flow]\label{def:NSContinuumCoherentSectorAndProjectedFlow}
Let
\[
  K_D
  \subset
  V_D
\]
denote the continuum coherent sector selected by the shared physical-sector gate and persistent-carrier extraction.
Let
\[
  \Theta_D(X)
\]
denote the coherent transport readout on that sector, and let
\[
  \mathcal L_{\mathrm{coh}}
\]
denote the dissipative coherent-sector generator.
The projected flow is the evolution on $K_D$ whose transport defect is measured against the coherent-sector equation
\[
  \partial_t \Theta_D(X)
  =
  - \nu \mathcal L_{\mathrm{coh}} \Theta_D(X)
  + E_{\mathrm{sh}}(X)
\]
with shadow error term $E_{\mathrm{sh}}(X)$.
\end{definition}

\begin{assumption}[NS projected-transport closure inputs]\label{as:NSProjectedTransportClosureInputs}
Assume the following packet on the continuum coherent sector.
\begin{enumerate}
  \item \textbf{Projected well-posedness.}
  The projected flow on $K_D$ is locally well posed on the required $H_D^s$ core.
  \item \textbf{Complement coercivity.}
  There exists $c_{\perp}>0$ such that the transport complement obeys the coercive estimate
  \[
    \langle \Delta_{D,\omega} Y, Y \rangle
    \ge
    c_{\perp} \|Y\|^2
  \]
  for all $Y \in K_D^{\perp}$ in the graph domain.
  \item \textbf{Lower-order leakage.}
  The off-carrier transport leakage is lower order relative to the coherent-sector energy and can be absorbed into the global continuation norm.
  \item \textbf{Exact shadow identity.}
  The stabilized readout kills the pure non-classical shadow defect so that the commutative readout of the projected coherent flow satisfies the classical transport law.
  \item \textbf{Continuation-level norm upgrade.}
  The coherent-sector energy estimate upgrades to the classical continuation norm required by the Clay endpoint.
\end{enumerate}
\end{assumption}

\begin{theorem}[Conditional projected transport]\label{thm:ConditionalNSProjectedTransportTheorem}
Under \cref{as:NSProjectedTransportClosureInputs}, the projected flow on the continuum coherent sector satisfies a closed transport equation modulo lower-order leakage, and the stabilized readout solves the classical Navier--Stokes equation on the same admissible time interval.
\end{theorem}

\begin{proof}
Projected well-posedness gives the local coherent-sector flow.
Complement coercivity and lower-order leakage keep the transport defect from generating new unstable support outside the coherent sector except through the controlled error term.
The exact shadow identity removes the pure non-classical remainder in the commutative readout.
The transported readout therefore satisfies the classical Navier--Stokes equation while the coherent-sector energy controls the same interval of existence.
\end{proof}

\begin{theorem}[NS classical closure packet]\label{thm:ConditionalNSClassicalClosurePacket}
Assume \cref{thm:ConditionalNSProjectedTransportTheorem} and the continuation-level norm upgrade from \cref{as:NSProjectedTransportClosureInputs}.
Then the classical closure burden is reduced to the two named frontier obligations:
\begin{enumerate}
  \item projected NC-flow well-posedness on the coherent sector, and
  \item the global coercive energy estimate that closes the continuation norm on the same theorem surface.
\end{enumerate}
If those two obligations are discharged, the projected-flow route eliminates the live singularity pathway on the classical surface.
\end{theorem}

\begin{proof}
The projected transport theorem identifies the coherent-sector route as the governing transport object.
Once projected well-posedness and the global coercive energy estimate both hold on that same route, the continuation norm closes and the classical blow-up mechanism is removed on the transported shadow surface.
\end{proof}

\begin{corollary}[NS import to the shared carrier picture]\label{cor:NSImportToTheSharedCarrierPicture}
Under the same assumptions, the Navier--Stokes route uses the shared selector, persistent invariant, and exact/persistent/tail split through the coherent-sector transport law rather than through an independent compactness package.
\end{corollary}

\begin{proof}
The coherent sector is selected by the shared physical-sector gate, the transport readout is carried by the same persistent layer used elsewhere in the appendix, and the exact and residual channels are controlled by the shared split assumptions.
The NS route therefore lands on the common carrier picture.
\end{proof}
